\subsection{欧几里得除法}

定理1. 设 \(a\) 与 \(b\) 是整数，使得 \(b > 0\) 。那么存在整数 \(q\) 与 \(r\) ，使得 \(a = bq + r\) 且 \(r < b\) ，且整数 \(q\) 与 \(r\) 由这些条件唯一确定。

关于 \(q\) 和 \(r\) 的条件等价于 \(bq \leqslant a < b(q + 1)\) 且 \(r = a - bq\) （第2号，命题2）。因此必须找到 \(q\) 使得 \(bq \leqslant a < b(q + 1)\) ；换言之， \(q\) 必须是使 \(a < b(q + 1)\) 成立的最小整数，这表明 \(q\) 和 \(r = a - bq\) 都唯一确定。为证明它们存在，注意到存在整数 \(p\) 使 \(a < bp\) ，例如 \(a + 1\) （因为 \(b > 0\) ）。设 \(m\) 为这些整数中最小的一个。由于 \(m \neq 0\) ，可写成 \(m = q + 1\) ，其中 \(q \leqslant m\) ( \(\S 4\) ，第2号，命题2）；于是 \(bq \leqslant a < b(q + 1)\) .

定义1. 采用定理1的记号，r称为a除以b的余数。如果r=0，我们说a是b的倍数，或a能被b整除，或b是a的除数，或b整除a，或b是a的因子。数q称为a除以b的商，记为 \(\frac{a}{b}\) 或a/b。

如果 \(a\) 不是 \(b\) 的倍数，数 \(q\) 称为 \(a\) 除以 \(b\) 的商的整数部分（参见一般拓扑学，第四章，第8节，第2号）。

在本章中，写作 \(a / b\) 或 \(\frac{a}{b}\) 将蕴含 \(b\) 整除 \(a\) .

关系 \(a = bq\) 并且 \(q = a / b\) 是等价的（如果 \(b > 0\) ）。每一个倍数 \(a'\) （它是倍数 \(a\) ——\(b\) 的一个倍数——的倍数）都是 \(b\) 的倍数，以及

\[
a ^ {\prime} / b = (a ^ {\prime} / a) (a / b) \quad \text { 若 } \quad a \neq 0.
\]

此外，如果 \(c\) 与 \(d\) 是 \(b\) 的倍数，那么 \(c + d\) 与 \(c - d\) （如果 \(d \leqslant c\) ）是 \(b\) 的倍数，并且我们有

\[
\frac {c + d}{b} = \frac {c}{b} + \frac {d}{b}, \quad \frac {c - d}{b} = \frac {c}{b} - \frac {d}{b}.
\]

是2的倍数的整数称为偶数，其余的称为奇数。根据定理1，奇数的形式为 \(2n + 1\) .

